% fig_geometry.tex -- M3: payload and rod geometry (TikZ). Dimensions from the rig (30 Sep 2026):
% plates at r 0.225 m on a 500 mm ring, 12 plates every 30 deg anticlockwise from +x, rod 0.55 m,
% pivot 4 cm below the UAV centre, working elevation 45 deg. Drawn to scale (1 m = 9 cm).
\begin{figure}[h]
\centering
\begin{tikzpicture}[x=9cm, y=9cm, >=Stealth, font=\small]
  % ---------------- (a) top view ----------------
  \begin{scope}
    \fill[gray!12] (0,0) circle (0.25);
    \fill[white] (0,0) circle (0.19);
    \draw[gray!60, thin] (0,0) circle (0.25) (0,0) circle (0.19);
    % the even triangle left after detaching plate 3
    \draw[densely dashed, gray!70] (30:0.225) -- (150:0.225) -- (270:0.225) -- cycle;
    % plates
    \foreach \k in {0,...,11} {
      \pgfmathsetmacro{\a}{30*\k}
      \node[draw, fill=white, minimum size=0.11cm, inner sep=0pt, rotate=\a] at (\a:0.225) {};
      \node[font=\scriptsize, text=gray!70!black] at (\a:0.29) {\k};
    }
    % the four carrying plates
    \foreach \k in {1,3,5,9} { \fill[black] (30*\k:0.225) circle (0.012); }
    \draw[thick, black] (90:0.225) circle (0.022);
    % centre, axes, rho
    \fill (0,0) circle (0.008);
    \draw[->, thin, gray] (0,0) -- (0.09,0) node[below, font=\scriptsize, text=black] {$x_L$};
    \draw[->, thin, gray] (0,0) -- (0,0.09) node[left, font=\scriptsize, text=black] {$y_L$};
    \draw[->, thick] (0,0) -- (30:0.213) node[pos=0.62, above left, inner sep=1pt] {$\boldsymbol{\rho}_i$};
    \node[font=\scriptsize, align=left, anchor=north] at (0,-0.33) {$\bullet$ carrying plates 1, 3, 5, 9\quad $\bigcirc$ plate 3, detached};
    \node[anchor=north west] at (-0.36,0.36) {(a)};
  \end{scope}
  % ---------------- (b) side view of one rod (A plate, P pivot, C UAV centre) ----------------
  \begin{scope}[shift={(0.62,-0.30)}]
    \fill[gray!12] (-0.06,0) rectangle (0.0,0.03);
    \draw[gray!60, thin] (-0.06,0) rectangle (0.0,0.03);
    \fill[black] (-0.012,0.03) rectangle (0.012,0.038);
    \coordinate (A) at (0,0.034);
    \coordinate (P) at (0.3889,0.4229);
    \coordinate (C) at (0.3889,0.4629);
    \draw[very thick] (A) -- (P);
    \fill (A) circle (0.007) (P) circle (0.007);
    \draw[thick] (0.2989,0.4629) -- (0.4789,0.4629);
    \draw[gray] (0.2989,0.4689) ellipse (0.035 and 0.006) (0.4789,0.4689) ellipse (0.035 and 0.006);
    \fill (C) circle (0.006);
    \draw[->, thick] (A) -- (0.1131,0.1471) node[pos=0.55, above left, inner sep=1pt] {$\mathbf{q}_i$};
    \draw[gray] (0.16,0.034) arc (0:45:0.16);
    \draw[gray, thin] (A) -- (0.22,0.034);
    \node[font=\scriptsize] at (0.185,0.1105) {45$^\circ$};
    \node[font=\scriptsize, anchor=north west] at (0.2045,0.2185) {$l_i = 0.55$\,m};
    \draw[gray, thin, densely dotted] (P) -- (0.5689,0.4229) (C) -- (0.5689,0.4629);
    \draw[<->, thin] (0.5689,0.4229) -- (0.5689,0.4629) node[midway, right, font=\scriptsize] {$\boldsymbol{\delta}$: 4\,cm};
    \node[font=\scriptsize, anchor=north] at (-0.03,-0.005) {ring};
    \node[font=\scriptsize, anchor=south] at (0.3889,0.49) {UAV $i$};
    \node[anchor=north west] at (-0.08,0.66) {(b)};
  \end{scope}
\end{tikzpicture}
\caption[Payload and rod geometry]{Payload and rod geometry, to scale. (a) Top view of the 0.86\,kg ring
with its twelve magnet plates at a radius of 0.225\,m, numbered anticlockwise from the $x_L$ axis. The
four-UAV layout uses plates 1, 3, 5 and 9; detaching the UAV on plate 3 leaves the even triangle
(dashed), under which the ring hangs level. (b) Side view of one rod at the working elevation: the rod
direction $\mathbf{q}_i$ points from the plate to the UAV, and the rod pivots $\boldsymbol{\delta}$ = 4\,cm
below the UAV's centre.}
\label{fig:geometry}
\end{figure}
